\scanpage{147}
\margin{\uline{parler de sous-schémas $\Sigma$ de $E\otimes K$.}}
\margin{\textcircled{\S\,2}}
Soient $E,K$ deux extensions de $\mathbb Q$, \deleted{$\Sigma$}
$\mathfrak U$ un idéal de \inserted{$E_K=$} $E\otimes_{\mathbb Q}K$,
$\overline K$ une clôture algébrique de $K$.
On dit que $\mathfrak U$ définit une CM-structure sur $E$, relativement à $K$,
si l'idéal $\mathfrak U'$ de
$E\otimes_{\mathbb Q}\unsure{\overline K^{\sharp}|\overline K}=E_K\otimes_K\overline K$
\uline{engendré} par $\mathfrak U$ correspond à un sous-ensemble \inserted{$\Sigma$}
de l'ensemble \deleted{$\unsure{I}$} \deleted{$I_E$}
$\operatorname{Spec}E_{\unsure{\overline K^{\sharp}|\overline K}}
\simeq\operatorname{Hom}_{\mathrm{alg}}(E,\overline K)$ qui est tel que pour tout
élément $\tau$ de $G=\operatorname{Gal}(\overline K/K)$ qui est une conjugaison,
on ait $\tau(\Sigma)=(\operatorname{Spec}E_{\overline K})-\Sigma$.
Noter que ceci implique que pour $g\in G$, on a aussi
$g\tau g^{-1}(\Sigma)=(\operatorname{Spec}E_{\overline K})-\Sigma=\tau\Sigma$,
donc $\tau^{-1}g\tau g^{-1}(\Sigma)=\Sigma$, donc $\Sigma$ est invariant
\deleted{dans} \inserted{par} le ss-groupe \deleted{\unsure{de}} fermé invariant $H$
engendré par les $\tau^{-1}g\tau g^{-1}$.
Cela signifie que, si $K_0\subset K$ est le corps de définition de l'idéal
$\mathfrak U$ de $E_K$, le corps $K_0$ \inserted{$=k_{\Sigma}$} est \uline{weilien};
\deleted{et} il est d'ailleurs weilien \uline{non réel}, car $\tau(\Sigma)\ne\Sigma$
[puisque sinon on aurait $\operatorname{Spec}E_{\overline K}=\varnothing$, ce qui n'est pas].

\deleted{Supposons maintenant également $K$ fini sur $\mathbb Q$.}
\inserted{Soit $\overline{\mathbb Q}$ une clôture \uline{algébrique} de $\mathbb Q$}
\deleted{Alors} de groupe de Galois $G$.
Alors $E,K$ correspondent à des espaces homogènes compacts sous $G$, soient $I_E,I_K$.
La donnée de $\mathfrak U$ correspond à celle d'une partie fermée de $I_E\times I_K$,
invariante par $G$. Choisissons un plongt $K\subset\overline{\mathbb Q}$,
donc $I_K\simeq G/Q$, $Q$ le sous-groupe de $G$ qui laisse fixe les éléments de $K$.

\deleted{L'hypothèse} \deleted{Remplaçant} \inserted{Prenant} $K$ pour \inserted{$\overline{\mathbb Q}=$}$\overline K$ a pour effet de remplacer
$I_K$ par son pt origine $y$ \deleted{et} (plongt de $K$ dans $\overline{\mathbb Q}$),
et $\Sigma$ devient $\Sigma_y=\Sigma\cap(I_E\times y)$.
\inserted{Soit $\Sigma'=(I_E\times I_K)-\Sigma$, d'où $\Sigma'_y$.}
L'hypothèse est que pour tte conj. complexe $\tau\in G$,
\[
\deleted{(\tau\times1)(\Sigma\cap(I_E\times y))=(I_E-\Sigma}
\qquad
(\tau\times1)\Sigma_y=\deleted{I_E-}\Sigma'_y
\]
\[
\deleted{(\tau\times1)\Sigma'_y=I_E-\Sigma'_y}
\]
i.e.
\par\noindent
\deleted{\mbox{\begin{minipage}[c]{\linewidth}
\unsure{que pour|que peut} écrire aussi
\[
\left\{\begin{array}{l}
(\tau\times1)\Sigma_y\subset I_E-\Sigma\\
(\tau\times1)\Sigma'_y\subset I_E-\unsure{\Sigma'_y|\Sigma_y}
\end{array}\right.
\]
\end{minipage}}}
\par
\margin{146}
